\documentclass[12pt, a4paper]{article}

% ===== Pacotes =====
\usepackage[utf8]{inputenc}
\usepackage[T1]{fontenc}
\usepackage[brazil]{babel}
\usepackage{geometry}
\usepackage{listings}
\usepackage{xcolor}
\usepackage{graphicx}
\usepackage{fancyhdr}
\usepackage{titlesec}
\usepackage{hyperref}
\usepackage{enumitem}
\usepackage{booktabs}
\usepackage{tabularx}
\usepackage{amsmath}

\geometry{left=2.5cm, right=2.5cm, top=3cm, bottom=3cm}

% ===== Cores para listings =====
\definecolor{codebg}{RGB}{245,245,245}
\definecolor{codegreen}{RGB}{0,128,0}
\definecolor{codegray}{RGB}{128,128,128}
\definecolor{codeblue}{RGB}{0,0,180}

% ===== Estilo de código Python =====
\lstdefinestyle{python}{
    language=Python,
    backgroundcolor=\color{codebg},
    basicstyle=\ttfamily\footnotesize,
    keywordstyle=\color{codeblue}\bfseries,
    stringstyle=\color{codegreen},
    commentstyle=\color{codegray}\itshape,
    numbers=left,
    numberstyle=\tiny\color{codegray},
    frame=single,
    breaklines=true,
    showstringspaces=false,
    tabsize=4,
    captionpos=b
}

% ===== Estilo de código C =====
\lstdefinestyle{clang}{
    language=C,
    backgroundcolor=\color{codebg},
    basicstyle=\ttfamily\footnotesize,
    keywordstyle=\color{codeblue}\bfseries,
    stringstyle=\color{codegreen},
    commentstyle=\color{codegray}\itshape,
    numbers=left,
    numberstyle=\tiny\color{codegray},
    frame=single,
    breaklines=true,
    showstringspaces=false,
    tabsize=4,
    captionpos=b
}

% ===== Estilo de código Prolog =====
\lstdefinestyle{prolog}{
    language=Prolog,
    backgroundcolor=\color{codebg},
    basicstyle=\ttfamily\footnotesize,
    keywordstyle=\color{codeblue}\bfseries,
    stringstyle=\color{codegreen},
    commentstyle=\color{codegray}\itshape,
    numbers=left,
    numberstyle=\tiny\color{codegray},
    frame=single,
    breaklines=true,
    showstringspaces=false,
    tabsize=4,
    captionpos=b
}

% ===== Estilo de saída =====
\lstdefinestyle{output}{
    backgroundcolor=\color{codebg},
    basicstyle=\ttfamily\scriptsize,
    frame=single,
    breaklines=true,
    showstringspaces=false
}

\begin{document}

% ===== CAPA =====
\begin{titlepage}
    \centering
    \vspace*{2cm}
    {\Large\textbf{Pontifícia Universidade Católica de Goiás}}\\[0.5cm]
    {\large Escola Politécnica e de Artes}\\[0.3cm]
    {\large Curso de Ciência da Computação}\\[3cm]
    {\LARGE\textbf{Atividade Final}}\\[0.5cm]
    {\Large Paradigmas de Linguagens de Programação}\\[0.3cm]
    {\large Sistema de Análise de Desempenho Acadêmico}\\[3cm]
    {\large\textbf{Alunos:}}\\[0.3cm]
    {\Large Antônio Calixto da Silva}\\
    {\large Amanda Barbosa Ferreira Agapito}\\
    {\large Guilherme José da Silva}\\
    {\large Matheus Silva Pains}\\[2cm]
    {\large\textbf{Professor:} Roussian de Ramos Alves Gaioso}\\[0.5cm]
    {\large Goiânia, Goiás}\\
    {\large 12 de Junho de 2026}
\end{titlepage}

\tableofcontents
\newpage

% ===================================================================
\section{Descrição do Problema}
% ===================================================================

O sistema recebe dados de alunos contendo nome, matrícula, três notas e frequência. A média é calculada como:

\[
\text{média} = \frac{\text{nota1} + \text{nota2} + \text{nota3}}{3}
\]

O aluno é \textbf{aprovado} se $\text{média} \geq 6{,}0$ \textbf{e} $\text{frequência} \geq 75\%$. Caso contrário, é \textbf{reprovado}.

\subsection{Dados de Entrada}

\begin{lstlisting}[style=output]
Ana;2024001;8.0;7.5;9.0;85
Bruno;2024002;5.0;6.0;4.5;80
Carlos;2024003;7.0;6.5;8.0;60
Daniela;2024004;9.0;9.5;10.0;95
Amanda;2024005;10;9;10;100
Guilherme;2024006;9;10;10;100
Matheus;2024007;10;10;9;100
\end{lstlisting}

\subsection{Resultado Esperado}

\begin{table}[h]
\centering
\begin{tabular}{llcccccl}
\toprule
\textbf{Nome} & \textbf{Matrícula} & \textbf{N1} & \textbf{N2} & \textbf{N3} & \textbf{Média} & \textbf{Freq} & \textbf{Situação} \\
\midrule
Ana       & 2024001 & 8,0  & 7,5  & 9,0  & 8,17 & 85  & Aprovado \\
Bruno     & 2024002 & 5,0  & 6,0  & 4,5  & 5,17 & 80  & Reprovado \\
Carlos    & 2024003 & 7,0  & 6,5  & 8,0  & 7,17 & 60  & Reprovado \\
Daniela   & 2024004 & 9,0  & 9,5  & 10,0 & 9,50 & 95  & Aprovado \\
Amanda    & 2024005 & 10,0 & 9,0  & 10,0 & 9,67 & 100 & Aprovado \\
Guilherme & 2024006 & 9,0  & 10,0 & 10,0 & 9,67 & 100 & Aprovado \\
Matheus   & 2024007 & 10,0 & 10,0 & 9,0  & 9,67 & 100 & Aprovado \\
Antonio   & 2024008 & 10,0 & 10,0 & 10,0 & 10,00 & 100 & Aprovado \\
\bottomrule
\end{tabular}
\end{table}

% ===================================================================
\section{Linguagens Escolhidas}
% ===================================================================

\begin{itemize}
    \item \textbf{C (imperativo)}
    \item \textbf{Python (orientado a objetos)}
    \item \textbf{Python (funcional)}
    \item \textbf{Prolog (lógico)}
\end{itemize}

C foi escolhida para a versão imperativa por ser a linguagem sugerida pelo enunciado e a mais representativa do paradigma procedural. Python foi utilizada nas versões OO e funcional por sua versatilidade. Prolog foi mantido por ser a referência natural do paradigma lógico.

% ===================================================================
\section{Parte 1 --- Implementação Imperativa (C)}
% ===================================================================

\lstinputlisting[style=clang, caption={parte1\_imperativo.c}]{parte1_imperativo.c}

\subsection{Saída da Execução}

\begin{lstlisting}[style=output]
Nome         Matricula    Media   Freq  Situacao
------------------------------------------------
Ana          2024001      8.17    85    Aprovado
Bruno        2024002      5.17    80    Reprovado
Carlos       2024003      7.17    60    Reprovado
Daniela      2024004      9.50    95    Aprovado
Amanda       2024005      9.67    100   Aprovado
Guilherme    2024006      9.67    100   Aprovado
Matheus      2024007      9.67    100   Aprovado
Antonio      2024008      10.00   100   Aprovado
------------------------------------------------
Total: 8 | Aprovados: 6 | Reprovados: 2
\end{lstlisting}

% ===================================================================
\section{Parte 2 --- Implementação Orientada a Objetos (Python)}
% ===================================================================

\lstinputlisting[style=python, caption={parte2\_oo.py}]{parte2_oo.py}

\subsection{Saída da Execução}

\begin{lstlisting}[style=output]
======================================================================
RELATORIO DE DESEMPENHO ACADEMICO - VERSAO OO (Python)
======================================================================
Nome         Matricula    N1     N2     N3     Media    Freq   Situacao
----------------------------------------------------------------------
Ana          2024001      8.0    7.5    9.0    8.17     85     Aprovado
Bruno        2024002      5.0    6.0    4.5    5.17     80     Reprovado
Carlos       2024003      7.0    6.5    8.0    7.17     60     Reprovado
Daniela      2024004      9.0    9.5    10.0   9.50     95     Aprovado
Amanda       2024005      10.0   9.0    10.0   9.67     100    Aprovado
Guilherme    2024006      9.0    10.0   10.0   9.67     100    Aprovado
Matheus      2024007      10.0   10.0   9.0    9.67     100    Aprovado
Antonio      2024008      10.0   10.0   10.0   10.00    100    Aprovado
----------------------------------------------------------------------
Total de alunos: 8
Aprovados: 6
Reprovados: 2
======================================================================
\end{lstlisting}

% ===================================================================
\section{Parte 3 --- Implementação Funcional (Python)}
% ===================================================================

\lstinputlisting[style=python, caption={parte3\_funcional.py}]{parte3_funcional.py}

\subsection{Saída da Execução}

\begin{lstlisting}[style=output]
======================================================================
RELATORIO DE DESEMPENHO ACADEMICO - VERSAO FUNCIONAL (Python)
======================================================================
Nome         Matricula    N1     N2     N3     Media    Freq   Situacao
----------------------------------------------------------------------
Ana          2024001      8.0    7.5    9.0    8.17     85     Aprovado
Bruno        2024002      5.0    6.0    4.5    5.17     80     Reprovado
Carlos       2024003      7.0    6.5    8.0    7.17     60     Reprovado
Daniela      2024004      9.0    9.5    10.0   9.50     95     Aprovado
Amanda       2024005      10.0   9.0    10.0   9.67     100    Aprovado
Guilherme    2024006      9.0    10.0   10.0   9.67     100    Aprovado
Matheus      2024007      10.0   10.0   9.0    9.67     100    Aprovado
Antonio      2024008      10.0   10.0   10.0   10.00    100    Aprovado
----------------------------------------------------------------------
Total de alunos: 8
Aprovados: 6
Reprovados: 2
======================================================================

--- Aprovados (via filter) ---
  Ana - Media: 8.17
  Daniela - Media: 9.50
  Amanda - Media: 9.67
  Guilherme - Media: 9.67
  Matheus - Media: 9.67
  Antonio - Media: 10.00
\end{lstlisting}

% ===================================================================
\section{Parte 4 --- Implementação Lógica (Prolog)}
% ===================================================================

\lstinputlisting[style=prolog, caption={parte4\_logico.pl}]{parte4_logico.pl}

\subsection{Exemplo de Consultas e Saída}

\begin{lstlisting}[style=output]
?- aprovado(X).
X = ana ;
X = daniela ;
X = amanda ;
X = guilherme ;
X = matheus ;
X = antonio.

?- reprovado(X).
X = bruno ;
X = carlos.

?- media(ana, M).
M = 8.166666666666666.
\end{lstlisting}

% ===================================================================
\section{Comparação de Legibilidade}
% ===================================================================

\begin{table}[h]
\centering
\begin{tabularx}{\textwidth}{lX}
\toprule
\textbf{Critério} & \textbf{Análise} \\
\midrule
Clareza dos nomes & Todas as versões Python usam nomes descritivos em português. Prolog usa átomos curtos (\texttt{ana}, \texttt{aprovado}), concisos mas claros. \\
\addlinespace
Organização & A versão OO é a mais organizada: dados e lógica ficam encapsulados em classes. A funcional separa bem entrada/processamento/saída. A imperativa é linear e direta. \\
\addlinespace
Detalhes sintáticos & Prolog tem a menor quantidade de sintaxe. Python imperativo exige mais linhas por usar \texttt{while} e índices explícitos. \\
\addlinespace
Fluxo do programa & A imperativa é a mais fácil de seguir passo a passo. A funcional exige entender composição de funções. Prolog não tem fluxo explícito. \\
\addlinespace
Localização das regras & Em Prolog, cada regra é uma cláusula isolada e fácil de encontrar. Na OO, estão nos métodos da classe. Na funcional, em funções nomeadas. \\
\bottomrule
\end{tabularx}
\caption{Comparação de legibilidade entre os paradigmas.}
\end{table}

\textbf{Conclusão:} Para leitores iniciantes, a versão imperativa é a mais legível. Para projetos maiores, a OO oferece melhor organização. Prolog é o mais conciso, mas exige familiaridade com o paradigma.

% ===================================================================
\section{Comparação de Facilidade de Escrita}
% ===================================================================

\begin{table}[h]
\centering
\begin{tabularx}{\textwidth}{lX}
\toprule
\textbf{Critério} & \textbf{Análise} \\
\midrule
Quantidade de código & Prolog: $\sim$40 linhas. Funcional: $\sim$70 linhas. Imperativa: $\sim$80 linhas. OO: $\sim$100 linhas. \\
\addlinespace
Dificuldade sintática & Python é uniforme nos três paradigmas. Prolog exige aprender unificação e backtracking. \\
\addlinespace
Manipulação de listas & A funcional é a mais natural com \texttt{map}/\texttt{filter}. A imperativa usa índices manuais. A OO itera com \texttt{for}. \\
\addlinespace
Estruturas auxiliares & A OO exige criar classes. A funcional usa dicionários. A imperativa usa listas de listas. Prolog usa fatos diretamente. \\
\addlinespace
Bibliotecas & Apenas a funcional importa \texttt{functools.reduce}. As demais usam apenas a biblioteca padrão. \\
\bottomrule
\end{tabularx}
\caption{Comparação de facilidade de escrita.}
\end{table}

\textbf{Conclusão:} A versão funcional Python foi a mais rápida de implementar pela expressividade de \texttt{map}/\texttt{filter}. Prolog é extremamente conciso, mas exige mudança de mentalidade.

% ===================================================================
\section{Comparação de Confiabilidade}
% ===================================================================

\begin{table}[h]
\centering
\begin{tabularx}{\textwidth}{lX}
\toprule
\textbf{Critério} & \textbf{Análise} \\
\midrule
Verificação de tipos & Python tem tipagem dinâmica --- erros de tipo só aparecem em tempo de execução. Prolog não tem tipos explícitos. \\
\addlinespace
Erros em tempo de execução & A imperativa é a mais suscetível a erros de índice (\texttt{alunos[i][5]}). A OO protege dados com encapsulamento. A funcional, ao usar funções puras, reduz efeitos colaterais. \\
\addlinespace
Entrada inválida & Nenhuma versão trata entradas inválidas de forma robusta. Seria necessário adicionar \texttt{try/except} em Python ou validação em Prolog. \\
\addlinespace
Acesso à memória & Python e Prolog não permitem acesso direto à memória, eliminando riscos de buffer overflow (diferente de C). \\
\addlinespace
Testabilidade & A funcional é a mais testável: funções puras sem efeito colateral podem ser testadas isoladamente. A OO permite testes unitários por método. \\
\bottomrule
\end{tabularx}
\caption{Comparação de confiabilidade.}
\end{table}

\textbf{Conclusão:} A versão funcional oferece maior confiabilidade por usar funções puras. A OO vem em segundo lugar pelo encapsulamento. Python e Prolog são mais seguros que C por não permitirem acesso direto à memória.

% ===================================================================
\section{Comparação sobre Nomes, Escopo e Tempo de Vida}
% ===================================================================

\subsection{Versão Imperativa}

\begin{itemize}
    \item \textbf{Nomes declarados:} \texttt{caminho\_arquivo}, \texttt{lista\_alunos} (globais); \texttt{nome}, \texttt{media}, \texttt{i} (locais em funções).
    \item \textbf{Visibilidade:} As variáveis globais são visíveis em todo o módulo. As locais existem apenas dentro de suas funções.
    \item \textbf{Escopo:} Funções como \texttt{calcular\_media()} têm escopo local; parâmetros e variáveis internas não vazam.
    \item \textbf{Tempo de vida:} Variáveis locais são criadas na chamada da função e destruídas ao retornar. As globais existem durante toda a execução.
\end{itemize}

\subsection{Versão Orientada a Objetos}

\begin{itemize}
    \item \textbf{Nomes declarados:} Classes \texttt{Aluno}, \texttt{Turma}; atributos \texttt{self.\_\_nome}, \texttt{self.\_\_nota1}; métodos \texttt{calcular\_media()}.
    \item \textbf{Visibilidade:} Atributos com \texttt{\_\_} (name mangling) são visíveis apenas dentro da classe. Métodos públicos são acessíveis externamente.
    \item \textbf{Escopo:} Cada objeto tem seu próprio escopo de atributos. Métodos acessam \texttt{self} para dados da instância.
    \item \textbf{Tempo de vida:} Objetos existem enquanto houver referência a eles. O garbage collector de Python os destrói quando não há mais referências.
\end{itemize}

\subsection{Versão Funcional}

\begin{itemize}
    \item \textbf{Nomes declarados:} Funções puras \texttt{calcular\_media()}, \texttt{verificar\_aprovacao()}, \texttt{gerar\_resultado()}; variável \texttt{alunos} no nível do módulo.
    \item \textbf{Visibilidade:} Cada função vê apenas seus parâmetros e variáveis locais. Não há estado mutável compartilhado.
    \item \textbf{Escopo:} Funções são de escopo global. Lambdas em \texttt{filter()} capturam variáveis do escopo envolvente (closure).
    \item \textbf{Tempo de vida:} Dados são criados a cada chamada e descartados após o retorno. Listas intermediárias (\texttt{resultados}) vivem no escopo do módulo.
\end{itemize}

\subsection{Versão Lógica (Prolog)}

\begin{itemize}
    \item \textbf{Nomes declarados:} Predicados \texttt{aluno/6}, \texttt{media/2}, \texttt{aprovado/1}; variáveis lógicas \texttt{Nome}, \texttt{Media}.
    \item \textbf{Visibilidade:} Fatos e regras são globais no programa. Variáveis lógicas existem apenas dentro de sua cláusula.
    \item \textbf{Escopo:} Cada cláusula é independente. A variável \texttt{Nome} em \texttt{aprovado(Nome)} só existe naquela regra.
    \item \textbf{Tempo de vida:} Variáveis são instanciadas durante a unificação e deixam de existir quando a consulta termina ou faz backtracking.
\end{itemize}

% ===================================================================
\section{Comparação sobre Tipos de Dados}
% ===================================================================

\begin{table}[h]
\centering
\begin{tabularx}{\textwidth}{lXX}
\toprule
\textbf{Aspecto} & \textbf{Python (3 versões)} & \textbf{Prolog} \\
\midrule
Tipagem & Dinâmica, forte & Sem tipos explícitos \\
\addlinespace
Tipos primitivos & \texttt{int}, \texttt{float}, \texttt{str}, \texttt{bool} & Átomos, números, variáveis \\
\addlinespace
Strings & Tipo nativo \texttt{str}, imutável & Átomos (\texttt{ana}, \texttt{bruno}) \\
\addlinespace
Coleções & Listas, dicionários, tuplas & Listas Prolog, termos compostos \\
\addlinespace
Objetos & Sim (versão OO) & Não existe o conceito \\
\addlinespace
Conversão de tipo & Explícita: \texttt{float()}, \texttt{int()} & Automática via \texttt{is} \\
\addlinespace
Erros de tipo & \texttt{TypeError} em tempo de execução & Falha na unificação \\
\bottomrule
\end{tabularx}
\caption{Comparação de tipos de dados.}
\end{table}

% ===================================================================
\section{Comparação entre os Paradigmas}
% ===================================================================

\begin{enumerate}
    \item \textbf{Como a versão imperativa organiza a solução?}\\
    Sequencialmente: lê dados, percorre com laço \texttt{while}, calcula média, verifica condição, imprime. O programador controla cada passo.

    \item \textbf{Como a versão orientada a objetos organiza a solução?}\\
    Modela o domínio: \texttt{Aluno} encapsula dados e comportamento, \texttt{Turma} gerencia a coleção. A lógica é distribuída em métodos.

    \item \textbf{Como a versão funcional organiza a solução?}\\
    Como pipeline de transformações: \texttt{ler\_dados $\rightarrow$ map(gerar\_resultado) $\rightarrow$ imprimir\_relatorio}. Dados fluem entre funções puras.

    \item \textbf{Como a versão declarativa organiza a solução?}\\
    Declara fatos (dados dos alunos) e regras (média, aprovação). O motor de inferência de Prolog encontra as respostas.

    \item \textbf{Qual versão ficou mais próxima da forma humana?}\\
    Prolog. A regra \texttt{aprovado(Nome) :- media(Nome, M), M >= 6.0, ...} é quase linguagem natural.

    \item \textbf{Qual versão ficou mais próxima da máquina?}\\
    A imperativa. O laço \texttt{while} com índice manual e atribuições explícitas reflete como o processador executa instruções.

    \item \textbf{Qual versão parece mais adequada para esse problema?}\\
    A funcional. O problema é essencialmente uma transformação de dados (entrada $\rightarrow$ processamento $\rightarrow$ saída), que é o caso de uso ideal do paradigma funcional.
\end{enumerate}

% ===================================================================
\section{Conclusão}
% ===================================================================

Cada paradigma demonstrou forças e limitações claras ao resolver o mesmo problema:

\begin{itemize}
    \item O \textbf{paradigma imperativo} é o mais intuitivo para iniciantes e oferece controle total sobre o fluxo, mas gera código mais verboso e propenso a erros de índice.

    \item O \textbf{paradigma orientado a objetos} brilha em organização e manutenibilidade. Encapsulamento protege os dados e facilita extensões futuras (ex.: adicionar novos tipos de aluno). É o mais adequado para sistemas grandes.

    \item O \textbf{paradigma funcional} produziu o código mais conciso e testável. Funções puras sem efeito colateral facilitam depuração e paralelismo. É ideal para problemas de transformação de dados.

    \item O \textbf{paradigma lógico} é o mais declarativo e conciso. O programador descreve \textit{o que} quer, não \textit{como} fazer. É poderoso para regras de negócio e sistemas especialistas, mas limitado para I/O e interfaces.
\end{itemize}

Não existe paradigma universalmente ``melhor''. A escolha depende do problema: sistemas corporativos favorecem OO; pipelines de dados favorecem funcional; regras de inferência favorecem lógico; scripts simples favorecem imperativo. O programador competente domina múltiplos paradigmas e escolhe o mais adequado para cada situação.

\end{document}
